\ifdefined\BookMaster
\def\BookEnd{}
\else
\documentclass[a4paper,12pt,fleqn,openany,twoside,fontset=windows]{ctexbook}

\usepackage[fleqn]{amsmath}
\usepackage{amssymb,amsthm,mathrsfs}
\usepackage{geometry}
\usepackage[dvipsnames]{xcolor}
\usepackage{graphicx}
\usepackage{tikz}
\usetikzlibrary{calc}
\usepackage[most]{tcolorbox}
\usepackage{fancyhdr}
\usepackage{titlesec}
\usepackage{tocloft}
\usepackage{enumitem}
\usepackage{microtype}
\usepackage{pifont}
\usepackage{hyperref}
\usepackage{romannum}
\usepackage{mathtools}

\definecolor{bookblue}{HTML}{264653}
\definecolor{bookteal}{HTML}{4F7C82}
\definecolor{booklight}{HTML}{EDF4F3}
\definecolor{bookgray}{HTML}{5E686B}

\geometry{
    inner=2.7cm,
    outer=2.5cm,
    top=2.7cm,
    bottom=2.7cm,
    headheight=18pt,
    headsep=18pt,
    footskip=25pt
}
\setlength{\mathindent}{0pt}
\setlength{\parindent}{5pt}
\setlength{\parskip}{3pt plus 1pt minus 1pt}
\linespread{1.25}
\allowdisplaybreaks[3]
\AtBeginDocument{%
    \setlength{\parindent}{5pt}
    \setlength{\abovedisplayskip}{8pt plus 2pt minus 2pt}
    \setlength{\belowdisplayskip}{8pt plus 2pt minus 2pt}
    \setlength{\abovedisplayshortskip}{5pt plus 1pt minus 1pt}
    \setlength{\belowdisplayshortskip}{5pt plus 1pt minus 1pt}
}

\hypersetup{
    unicode=true,
    colorlinks=true,
    linkcolor=bookblue,
    citecolor=bookblue,
    urlcolor=bookteal,
    bookmarksopen=true,
    pdfauthor={Chen},
    pdftitle={数学分析习题整理}
}

\renewcommand{\chaptermark}[1]{%
    \edef\@tmp{第\thechapter 章\quad #1}%
    \expandafter\markboth\expandafter{\@tmp}{}%
}
\fancypagestyle{main}{%
    \fancyhf{}
    \fancyhead[LE]{\small\color{bookgray}\nouppercase{\leftmark}}
    \fancyhead[RO]{\small\color{bookgray}数学分析习题整理}
    \fancyfoot[C]{\color{bookblue}\thepage}
    \fancyfoot[R]{\small\bfseries Chen\;\; github.com/mathlover-1/math-analysis}
    \renewcommand{\headrulewidth}{0.4pt}
    \renewcommand{\footrulewidth}{0pt}
}
\fancypagestyle{plain}{%
    \fancyhf{}
    \fancyfoot[C]{\color{bookblue}\thepage}
    \fancyfoot[R]{\small\bfseries Chen\;\; github.com/mathlover-1/math-analysis}
    \renewcommand{\headrulewidth}{0pt}
    \renewcommand{\footrulewidth}{0pt}
}
\pagestyle{main}

\titleformat{\chapter}[display]
  {\normalfont\sffamily\bfseries\color{bookblue}}
  {\raggedleft\fontsize{46}{50}\selectfont\thechapter}
  {-18pt}
  {\titlerule[1.2pt]\vspace{10pt}\filright\Huge}
\titlespacing*{\chapter}{0pt}{-15pt}{36pt}

\renewcommand{\contentsname}{目\quad 录}
\renewcommand{\cfttoctitlefont}{\hfill\sffamily\bfseries\Huge\color{bookblue}}
\renewcommand{\cftaftertoctitle}{\hfill}
\setlength{\cftbeforetoctitleskip}{5pt}
\setlength{\cftaftertoctitleskip}{28pt}
\setlength{\cftbeforechapskip}{12pt}
\renewcommand{\cftchapfont}{\sffamily\bfseries\large\color{bookblue}}
\renewcommand{\cftchappagefont}{\sffamily\bfseries\large\color{bookblue}}
\renewcommand{\cftchapleader}{\color{bookteal}\cftdotfill{2.5}}

\newtcolorbox{ti}{
    enhanced,
    breakable,
    colback=booklight!55!white,
    colframe=bookteal!35!white,
    boxrule=0.45pt,
    boxsep=0pt,
    left=10pt,
    right=10pt,
    top=9pt,
    bottom=9pt,
    before skip=14pt,
    after skip=14pt,
    arc=2pt,
    outer arc=2pt,
    before upper={\parindent=0pt}
}

\newenvironment{booknotebody}{%
    \begin{center}
    \begin{minipage}{0.84\textwidth}
    \songti\zihao{-4}\linespread{1.45}\selectfont
    \setlength{\parindent}{2\ccwd}
    \setlength{\parskip}{0.75em}
}{%
    \end{minipage}
    \end{center}
}

\newenvironment{booknotelongbody}{%
    \begingroup
    \songti\zihao{-4}\linespread{1.45}\selectfont
    \setlength{\parindent}{2\ccwd}
    \setlength{\parskip}{0.75em}
    \begin{list}{}{%
        \setlength{\leftmargin}{0.08\textwidth}
        \setlength{\rightmargin}{0.08\textwidth}
        \setlength{\topsep}{0pt}
        \setlength{\partopsep}{0pt}
        \setlength{\parsep}{0pt}
        \setlength{\itemsep}{0pt}
    }
    \item\relax
}{%
    \end{list}
    \endgroup
}

\fancypagestyle{plainnowm}{%
    \fancyhf{}
    \fancyfoot[C]{\color{bookblue}\thepage}
    \renewcommand{\headrulewidth}{0pt}
    \renewcommand{\footrulewidth}{0pt}
}

\newcommand{\booknotetitle}[2][plain]{%
    \clearpage
    \phantomsection
    \addcontentsline{toc}{chapter}{#2}
    \markboth{#2}{}
    \thispagestyle{#1}
    \vspace*{1.2cm}
    \begin{center}
        {\sffamily\heiti\bfseries\fontsize{26}{34}\selectfont\color{bookblue}#2\par}
        \vspace{0.8em}
        {\color{bookteal}\rule{4.5cm}{0.8pt}\par}
    \end{center}
    \vspace{0.6cm}
}

\renewcommand{\proofname}{证明}
\renewcommand{\qedsymbol}{\color{bookteal}\ensuremath{\blacksquare}}

\title{数学分析习题整理}
\author{Chen}
\date{}

\makeatletter
% ==================== 旧版封面/封底设计备份 ====================
% 2026-08 重设计,旧代码用 \iffalse...\fi 封存,新版见下方。
\iffalse
\renewcommand{\maketitle}{%
    \begin{titlepage}
        \thispagestyle{empty}
        \setcounter{page}{0}
        \noindent\hspace*{-\parindent}\begin{tikzpicture}[remember picture,overlay]
            \fill[bookblue] (current page.north west) rectangle ([yshift=-4.2cm]current page.north east);
            \fill[booklight] ([yshift=-4.2cm]current page.north west) rectangle (current page.south east);
            \draw[bookteal!55,line width=0.7pt]
                ([xshift=2.4cm,yshift=4.3cm]current page.south west)
                .. controls +(2.5cm,1.8cm) and +(-2.8cm,-1.3cm) ..
                ([xshift=-2.4cm,yshift=7.2cm]current page.south east);
            \draw[bookteal!35,line width=0.7pt]
                ([xshift=2.5cm,yshift=6.2cm]current page.south west)
                .. controls +(3.0cm,-1.2cm) and +(-3.0cm,1.6cm) ..
                ([xshift=-2.3cm,yshift=3.8cm]current page.south east);
        \end{tikzpicture}
        \vspace*{5.6cm}
        \noindent\color{bookblue}\rule{\textwidth}{1.4pt}\par
        \vspace{8pt}
        {\centering\sffamily\bfseries\fontsize{32}{42}\selectfont\color{bookblue}\@title\par}
        \vspace{8pt}
        \noindent\color{bookteal}\rule{\textwidth}{0.6pt}\par
        \vspace{2.6cm}
        {\centering\Large\color{bookgray}\@author\par}
        \vfill
        {\centering\color{bookteal}\large
            $\displaystyle \lim_{n\to\infty} a_n$\qquad
            $\displaystyle \int_a^b f(x)\,\mathrm{d}x$\qquad
            $\displaystyle \nabla f(x,y)$\par
        }
        \vspace*{1.5cm}
    \end{titlepage}
}

\newcommand{\makebackcover}{%
    \begingroup
    \let\cleardoublepage\clearpage
    \begin{titlepage}
        \thispagestyle{empty}
        \noindent\hspace*{-\parindent}\begin{tikzpicture}[remember picture,overlay]
            \fill[bookblue] (current page.north west) rectangle (current page.south east);
            \fill[booklight,rounded corners=7pt]
                ([xshift=2.2cm,yshift=-3.0cm]current page.north west)
                rectangle
                ([xshift=-2.2cm,yshift=3.0cm]current page.south east);
            \fill[bookteal!80]
                (current page.north west) --
                ([xshift=5.3cm]current page.north west) --
                ([xshift=2.1cm,yshift=-4.4cm]current page.north west) --
                ([yshift=-6.2cm]current page.north west) -- cycle;
            \fill[bookteal!55]
                (current page.south east) --
                ([xshift=-5.8cm]current page.south east) --
                ([xshift=-2.4cm,yshift=4.2cm]current page.south east) --
                ([yshift=6.0cm]current page.south east) -- cycle;
            \foreach \radius in {0.75,1.15,1.55,1.95} {
                \draw[bookteal!55,line width=0.55pt]
                    ([xshift=-2.9cm,yshift=-4.4cm]current page.north east)
                    circle[radius=\radius cm];
            }
            \foreach \radius in {0.65,1.05,1.45} {
                \draw[booklight!65,line width=0.55pt]
                    ([xshift=2.4cm,yshift=3.6cm]current page.south west)
                    circle[radius=\radius cm];
            }
            \fill[bookteal!16,rounded corners=3pt]
                ([xshift=3.5cm,yshift=-7.1cm]current page.north west)
                rectangle
                ([xshift=-5.0cm,yshift=7.2cm]current page.south east);
            \fill[bookteal!30,rounded corners=3pt]
                ([xshift=4.4cm,yshift=-8.0cm]current page.north west)
                rectangle
                ([xshift=-4.1cm,yshift=8.1cm]current page.south east);
            \begin{scope}[shift={([xshift=-0.75cm]current page.center)}]
                \clip[rounded corners=3pt] (-5.25,-4.75) rectangle (5.25,4.75);
                \foreach \v in {-3.6,-3.0,...,3.6} {
                    \draw[bookteal!48,line width=0.45pt,opacity=0.72]
                        plot[domain=-3.6:3.6,samples=45,smooth,variable=\u]
                        ({\u+0.55*\v},{0.28*\u-0.35*\v+0.11*(\u*\u-\v*\v)});
                }
                \foreach \u in {-3.6,-3.0,...,3.6} {
                    \draw[bookblue!42,line width=0.42pt,opacity=0.62]
                        plot[domain=-3.6:3.6,samples=45,smooth,variable=\v]
                        ({\u+0.55*\v},{0.28*\u-0.35*\v+0.11*(\u*\u-\v*\v)});
                }
                \draw[bookteal!70,line width=0.75pt,opacity=0.82]
                    plot[domain=-3.6:3.6,samples=55,smooth,variable=\u]
                    ({\u},{0.28*\u+0.11*\u*\u});
                \draw[bookblue!60,line width=0.7pt,opacity=0.75]
                    plot[domain=-3.6:3.6,samples=55,smooth,variable=\v]
                    ({0.55*\v},{-0.35*\v-0.11*\v*\v});
            \end{scope}
            \draw[bookteal!45,line width=0.65pt]
                ([xshift=3.0cm,yshift=-8.2cm]current page.north west)
                .. controls +(2.2cm,2.0cm) and +(-2.0cm,-2.1cm) ..
                ([xshift=-3.0cm,yshift=9.0cm]current page.south east);
            \draw[bookteal!25,line width=0.55pt]
                ([xshift=3.1cm,yshift=8.0cm]current page.south west)
                .. controls +(2.7cm,-1.7cm) and +(-2.6cm,1.8cm) ..
                ([xshift=-3.1cm,yshift=-8.7cm]current page.north east);
            \node[anchor=west,text=bookteal!72,scale=0.78] at
                ([xshift=2.8cm,yshift=-4.0cm]current page.north west) {%
                $\displaystyle e^{\mathrm{i}\pi}+1=0$};
            \node[anchor=east,text=bookblue!62,scale=0.67] at
                ([xshift=-2.7cm,yshift=4.0cm]current page.south east) {%
                $\displaystyle
                f(x)=\sum_{n=0}^{\infty}\frac{f^{(n)}(a)}{n!}(x-a)^n$};
        \end{tikzpicture}
        \mbox{}
    \end{titlepage}
    \endgroup
}
\fi

% ==================== 新版封面(2026-08 重设计) ====================
% 设计要点:顶部深绿纵向渐变+波浪下缘+两层半透明过渡波(解决深浅绿生硬过渡);
% 深色区螺旋线/同心圆母题;标题+英文副题+作者居中;
% 中部 Fourier 部分和逼近方波曲线;底部浅色波纹(无文字)。
\renewcommand{\maketitle}{%
    \begin{titlepage}
        \thispagestyle{empty}
        \setcounter{page}{0}
        \noindent\hspace*{-\parindent}\begin{tikzpicture}[remember picture,overlay]
            % ---------- 底色 ----------
            \fill[booklight] (current page.north west) rectangle (current page.south east);

            % ---------- 顶部深色区:纵向渐变 + 波浪下缘 ----------
            \shade[top color=bookblue, bottom color=bookteal!55!bookblue]
                (current page.north west) -- (current page.north east) --
                ([yshift=-8.6cm]current page.north east)
                .. controls ([xshift=-5.2cm,yshift=-10.6cm]current page.north east)
                         and ([xshift=6.0cm,yshift=-8.0cm]current page.north west) ..
                ([yshift=-10.0cm]current page.north west) -- cycle;

            % ---------- 过渡层 1:半透明波浪 ----------
            \fill[bookteal!60!bookblue,opacity=0.42]
                ([yshift=-8.6cm]current page.north east)
                .. controls ([xshift=-5.2cm,yshift=-10.6cm]current page.north east)
                         and ([xshift=6.0cm,yshift=-8.0cm]current page.north west) ..
                ([yshift=-10.0cm]current page.north west) --
                ([yshift=-12.6cm]current page.north west)
                .. controls ([xshift=6.5cm,yshift=-11.0cm]current page.north west)
                         and ([xshift=-5.5cm,yshift=-13.2cm]current page.north east) ..
                ([yshift=-11.4cm]current page.north east) -- cycle;

            % ---------- 过渡层 2:更浅的波浪 ----------
            \fill[bookteal!35,opacity=0.35]
                ([yshift=-10.6cm]current page.north east)
                .. controls ([xshift=-6.0cm,yshift=-12.4cm]current page.north east)
                         and ([xshift=5.0cm,yshift=-9.8cm]current page.north west) ..
                ([yshift=-11.8cm]current page.north west) --
                ([yshift=-13.8cm]current page.north west)
                .. controls ([xshift=5.5cm,yshift=-12.6cm]current page.north west)
                         and ([xshift=-6.5cm,yshift=-14.6cm]current page.north east) ..
                ([yshift=-12.8cm]current page.north east) -- cycle;

            % ---------- 深色区装饰:右上螺旋线 ----------
            \begin{scope}[shift={([xshift=-3.4cm,yshift=-4.3cm]current page.north east)}]
                \draw[booklight!45!bookteal,opacity=0.55,line width=0.6pt]
                    plot[domain=0:1440,samples=220,smooth,variable=\t]
                    ({0.0028*\t*cos(\t)},{0.0028*\t*sin(\t)});
            \end{scope}

            % ---------- 深色区装饰:左下同心圆 ----------
            \foreach \r in {0.7,1.25,1.8,2.35,2.9}{
                \draw[booklight!35!bookteal,opacity=0.5,line width=0.5pt]
                    ([xshift=0.4cm,yshift=-7.9cm]current page.north west) circle[radius=\r cm];
            }

            % ---------- 主标题 ----------
            \node[anchor=center,text=booklight] at
                ([yshift=-4.5cm]current page.north)
                {\sffamily\bfseries\fontsize{40}{46}\selectfont \@title};

            % ---------- 分隔短线 ----------
            \draw[booklight!65,line width=1.2pt]
                ([xshift=-2.6cm,yshift=-5.75cm]current page.north) --
                ([xshift=2.6cm,yshift=-5.75cm]current page.north);

            % ---------- 英文副题 ----------
            \node[anchor=center,text=booklight!70!bookteal] at
                ([yshift=-6.55cm]current page.north)
                {\itshape\large Mathematical Analysis \textperiodcentered\
                 A Collection of Exercises};

            % ---------- 作者(紧随标题下方) ----------
            \node[anchor=center,text=booklight!88] at
                ([yshift=-8.2cm]current page.north)
                {\sffamily\large \@author\quad 整理};

            % ---------- 中部:Fourier 部分和逼近方波 ----------
            \begin{scope}[shift={([xshift=-5.2cm,yshift=-18.1cm]current page.north)}]
                % 坐标轴
                \draw[bookgray!55,line width=0.5pt,->] (-0.5,0) -- (10.9,0);
                \draw[bookgray!55,line width=0.5pt,->] (0,-1.75) -- (0,1.95);
                % 目标方波(虚线)
                \draw[bookgray!75,line width=0.6pt,dashed]
                    (0,0.92) -- (5,0.92) (5,-0.92) -- (10,-0.92);
                % N=1
                \draw[bookteal!45,line width=0.7pt]
                    plot[domain=0:360,samples=90,smooth,variable=\x]
                    ({\x/36},{1.18*sin(\x)});
                % N=3
                \draw[bookteal!80,line width=0.8pt]
                    plot[domain=0:360,samples=120,smooth,variable=\x]
                    ({\x/36},{1.18*(sin(\x)+sin(3*\x)/3)});
                % N=7
                \draw[bookblue,line width=0.95pt]
                    plot[domain=0:360,samples=150,smooth,variable=\x]
                    ({\x/36},{1.18*(sin(\x)+sin(3*\x)/3+sin(5*\x)/5+sin(7*\x)/7)});
                % 刻度标注
                \node[anchor=north east,text=bookgray!80,scale=0.62] at (-0.12,0) {$0$};
                \node[anchor=north,text=bookgray!80,scale=0.62] at (5,0) {$\pi$};
                \node[anchor=north,text=bookgray!80,scale=0.62] at (10,0) {$2\pi$};
            \end{scope}
            % 图注
            \node[anchor=center,text=bookgray,scale=0.8] at
                ([yshift=-20.5cm]current page.north)
                {$S_N(x)=\displaystyle\sum_{k=1}^{N}\frac{\sin\bigl((2k-1)x\bigr)}{2k-1}$
                 \ 逐步逼近方波};

            % ---------- 底部浅色波纹(无文字,与顶部波浪呼应) ----------
            \fill[bookteal!22,opacity=0.5]
                ([yshift=2.4cm]current page.south west)
                .. controls ([xshift=6.0cm,yshift=3.9cm]current page.south west)
                         and ([xshift=-6.5cm,yshift=1.0cm]current page.south east) ..
                ([yshift=2.9cm]current page.south east) --
                (current page.south east) -- (current page.south west) -- cycle;
            \fill[bookteal!40!bookblue,opacity=0.28]
                ([yshift=1.3cm]current page.south west)
                .. controls ([xshift=5.5cm,yshift=2.7cm]current page.south west)
                         and ([xshift=-5.5cm,yshift=0.3cm]current page.south east) ..
                ([yshift=1.9cm]current page.south east) --
                (current page.south east) -- (current page.south west) -- cycle;
        \end{tikzpicture}
        \mbox{}
    \end{titlepage}
}

% ==================== 新版封底(纯背景,2026-08 重设计) ====================
% 设计要点:对角渐变底+角部半透明叠层;中下部外摆线(玫瑰形)为视觉主体;
% 底部两层半透明白色波浪;无任何文字。
\newcommand{\makebackcover}{%
    \begingroup
    \let\cleardoublepage\clearpage
    \begin{titlepage}
        \thispagestyle{empty}
        \noindent\hspace*{-\parindent}\begin{tikzpicture}[remember picture,overlay]
            % ---------- 底色:对角渐变 ----------
            \shade[top color=bookblue, bottom color=bookteal!48!bookblue,
                   shading angle=-30]
                (current page.north west) rectangle (current page.south east);

            % ---------- 角部半透明叠层 ----------
            \fill[white,opacity=0.05]
                ([xshift=2.0cm,yshift=1.0cm]current page.north east) circle[radius=7.2cm];
            \fill[white,opacity=0.06]
                ([xshift=-0.5cm,yshift=-1.5cm]current page.north east) circle[radius=3.6cm];

            % ---------- 左下同心圆 ----------
            \foreach \r in {1.0,1.75,2.5,3.25,4.0,4.75}{
                \draw[booklight!40,opacity=0.55,line width=0.55pt]
                    ([xshift=-0.6cm,yshift=-0.4cm]current page.south west) circle[radius=\r cm];
            }
            % ---------- 右下螺旋线 ----------
            \begin{scope}[shift={([xshift=-2.9cm,yshift=3.2cm]current page.south east)}]
                \draw[booklight!45!bookteal,opacity=0.5,line width=0.55pt]
                    plot[domain=0:1260,samples=200,smooth,variable=\t]
                    ({0.0036*\t*cos(\t)},{0.0036*\t*sin(\t)});
            \end{scope}
            % ---------- 左上斜线组 ----------
            \foreach \i in {0,1,2}{
                \draw[bookteal!60!booklight,opacity=0.5,line width=0.6pt]
                    ([xshift={1.6cm+\i*0.45cm}]current page.north west) --
                    ([yshift={-1.6cm-\i*0.45cm}]current page.north west);
            }

            % ---------- 中下部视觉主体:外摆线(玫瑰形) ----------
            \begin{scope}[shift={([yshift=-16.6cm]current page.north)}]
                % 外圈
                \draw[booklight!35,opacity=0.6,line width=0.6pt]
                    (0,0) circle[radius=4.9cm];
                \draw[booklight!22,opacity=0.5,line width=0.5pt]
                    (0,0) circle[radius=5.25cm];
                % 主玫瑰线
                \draw[booklight!55!bookteal,opacity=0.8,line width=0.85pt]
                    plot[domain=0:1080,samples=400,smooth,variable=\t]
                    ({0.335*(8*cos(\t)-5*cos(2.6667*\t))},
                     {0.335*(8*sin(\t)-5*sin(2.6667*\t))});
                % 内层玫瑰线
                \draw[booklight!38!bookteal,opacity=0.6,line width=0.6pt]
                    plot[domain=0:1080,samples=400,smooth,variable=\t]
                    ({0.21*(8*cos(\t)-5*cos(2.6667*\t))},
                     {0.21*(8*sin(\t)-5*sin(2.6667*\t))});
            \end{scope}

            % ---------- 底部半透明波浪 ----------
            \fill[white,opacity=0.06]
                ([yshift=4.6cm]current page.south west)
                .. controls ([xshift=6.5cm,yshift=6.4cm]current page.south west)
                         and ([xshift=-6.0cm,yshift=2.6cm]current page.south east) ..
                ([yshift=5.2cm]current page.south east) --
                (current page.south east) -- (current page.south west) -- cycle;
            \fill[booklight,opacity=0.08]
                ([yshift=2.6cm]current page.south west)
                .. controls ([xshift=5.5cm,yshift=4.3cm]current page.south west)
                         and ([xshift=-6.5cm,yshift=1.0cm]current page.south east) ..
                ([yshift=3.3cm]current page.south east) --
                (current page.south east) -- (current page.south west) -- cycle;
        \end{tikzpicture}
        \mbox{}
    \end{titlepage}
    \endgroup
}
\makeatother

% 页脚右下方常态化水印(署名 + 仓库地址,加粗黑色),已在 main/plain 分页样式中定义,所有页面均生效.

\begin{document}
\mainmatter
\pagestyle{main}
\setcounter{chapter}{7}
\def\BookEnd{\end{document}}
\fi

% ==================== 附录:数学分析重要定理 ====================
% 本附录按小节组织,小节编号 §1,§2,§3,... 自动生成,
% 你只需写 \section{小节名},例如 \section{数列}.
% 子标题自动居左,编号前带段落符号 \S.
%
% 定理请放在 ti 环境(绿色框)中,证明写在框外下方,与正文题目结构一致:
%   \begin{ti}
%   【定理1】 定理陈述...
%   \end{ti}
%   \textbf{证明}:证明内容...
%
% 同样支持 \textbf{另解}、\textbf{注}、\textbf{注记}、\textbf{一点思考} 等标记.
% ================================================================

% 子标题编号格式:§1,§2,§3,...(带段落符号);标题居左.
\renewcommand{\thesection}{\S\arabic{section}}
\setcounter{section}{0}
\titleformat{\section}[block]
  {\normalfont\heiti\bfseries\Large\color{bookblue}}
  {\thesection}{1em}{}
\titlespacing*{\section}{0pt}{3.2em plus 0.4em}{1em}

% 子节(subsection)编号格式:§X.Y(父节号.子节号),标题居左,比小节略小.
\renewcommand{\thesubsection}{\S\arabic{section}.\arabic{subsection}}
\titleformat{\subsection}[block]
  {\normalfont\heiti\bfseries\large\color{bookblue}}
  {\thesubsection}{1em}{}
\titlespacing*{\subsection}{0pt}{2.4em plus 0.3em}{0.8em}

% 附录行间编号公式居中:正文已按 fleqn(左对齐)排版,本附录要求 equation/align 等编号公式居中,
% 故在此关闭 amsmath 的 fleqn 布局分支(无编号 $$...$$ 公式本就居中,不受影响)。仅作用于本附录及其后内容。
\makeatletter\@fleqnfalse\makeatother

% 附录公式编号:去掉"章.节"前缀,同一节内从 1 开始递推(1,2,3,...)。
\makeatletter
\@addtoreset{equation}{section}
\renewcommand{\theequation}{\arabic{equation}}
\makeatother

% 附录内局部放大绿框(ti)之间的间距:+20%(14pt → 16.8pt),不影响正文各章.
\renewtcolorbox{ti}{
    enhanced,
    breakable,
    colback=booklight!55!white,
    colframe=bookteal!35!white,
    boxrule=0.45pt,
    boxsep=0pt,
    left=10pt,
    right=10pt,
    top=9pt,
    bottom=9pt,
    before skip=16.8pt,
    after skip=16.8pt,
    arc=2pt,
    outer arc=2pt,
    before upper={\parindent=0pt}
}

\setcounter{section}{4}   % 从§5开始编号
% ==================== 附录:数学分析重要定理(第五至六节) ====================
% 本文件为附录的中段:§5 一元微分学与一元积分学 与 §6 多元微分学与多元积分学.
% 前段(§1-§4)见 附录-数学分析重要定理-第1至4节.tex.
% 三文件均可独立编译;主文件中按顺序 input 三者即可.
% ========================================================================
\section{一元微分学与一元积分学}
\subsection{一元函数的微分}
\begin{ti}
\textbf{定义5.1(导数)}  设函数$f$在$x_0$附近有定义.如果极限
$$\lim_{x\to x_0}\frac{f(x)-f(x_0)}{x-x_0}$$
存在且有限,则称$f$在$x_0$处\textbf{可导},此极限称为$f$在$x_0$处的\textbf{导数},记为$f'(x_0)$.
\end{ti}


\begin{ti}
\textbf{命题5.2(导数的四则运算)}  设$f,g$在$x_0$处可导,$\lambda,\mu\in\mathbb{R}$,则$fg$和$\lambda f+\mu g$也在$x_0$处可导,且\\
(1)$(\lambda f+\mu g)'(x_0)=\lambda f'(x_0)+\mu g'(x_0)$;\\
(2)$(fg)'(x_0)=f'(x_0)g(x_0)+f(x_0)g'(x_0)$;\\
(3)当$g(x_0)\neq 0$时,$\dfrac{f}{g}$也在$x_0$处可导,且$\left(\dfrac{f}{g}\right)'(x_0)=\dfrac{f'(x_0)g(x_0)-f(x_0)g'(x_0)}{g^2(x_0)}$.
\end{ti}

\begin{ti}
\textbf{定理5.3(链式法则)}  设$g$在$x_0$处可导,$f$在$g(x_0)$处可导,则复合函数$f\circ g=f(g)$在$x_0$处可导,且
$$[f(g)]'(x_0)=f'(g(x_0))g'(x_0)$$.
\end{ti}

\begin{ti}
\textbf{定理5.4(反函数求导法则)}  设$f$在$x_0$附近连续且有反函数$g$.如果$f$在$x_0$处可导,且导数$f'(x_0)\neq 0$,则$g$在$y_0=f(x_0)$处可导,且
$$g'(y_0)=\left[f'(x_0)\right]^{-1}$$.
\end{ti}

\begin{ti}
\textbf{定义5.5(微分)}  设函数$f$在$x_0$的某邻域中有定义.如果存在$\alpha\in\mathbb{R}$,使得
$$f(x)=f(x_0)+\alpha(x-x_0)+o(x-x_0)\,(x\rightarrow x_0),$$
则称$f$在$x_0$处\textbf{可微},线性映射
$$df(x_0):\mathbb{R} \rightarrow \mathbb{R} ,t\rightarrow \alpha t$$
称为$f$在$x_0$处的\textbf{微分}.
\end{ti}

\begin{ti}
\textbf{定义5.6(极值)}  设函数$f$在$x_0$的某邻域中有定义.如果存在$\delta>0$,使得当$\left\lvert x-x_0\right\rvert<\delta$时$f(x)\leq f(x_0)$恒成立,
则称$x_0$为$f$的\textbf{极大值点},$f(x_0)$为\textbf{极大值};极小值点与极小值类似定义.极大值与极小值统称\textbf{极值}.
\end{ti}

\begin{ti}
\textbf{定理5.7($Fermat$引理)}  设$x_0$为函数$f$在$I$中的极值点.如果$f$在$x_0$处可导,且$x_0$为$I$的内点,则$f'(x_0)=0$.
\end{ti}

\textbf{证明}:不妨设$x_0$为$f$的极大值点(不然可考虑$-f$).由于$x_0$为$I$的内点,故存在$\delta>0$,使得$(x_0-\delta,x_0+\delta)\subset I$.取$\delta$充分小,使得$f(x)\leq f(x_0)$对一切$x\in(x_0-\delta,x_0+\delta)$成立.特别地,当$x\in(x_0-\delta,x_0)$时,
$$f'(x_0)=f'_-(x_0)=\lim_{x\to x_0^-}\frac{f(x)-f(x_0)}{x-x_0}\geq 0;$$
\phantom{\textbf{证明}:}而当$x\in(x_0,x_0+\delta)$时,
$$f'(x_0)=f'_+(x_0)=\lim_{x\to x_0^+}\frac{f(x)-f(x_0)}{x-x_0}\leq 0.$$
这说明$f'(x_0)=0$.

\begin{ti}
\textbf{定理5.8($Rolle$中值定理)}  设函数$f$在$[a,b]$上连续,在$(a,b)$内可导,且$f(a)=f(b)$,则存在$\xi\in(a,b)$,使得$f'(\xi)=0$.
\end{ti}

\textbf{证明}:根据最值定理,$f$有最大值和最小值.如果二者相等,则$f$为常值函数,定理自然成立;否则由$f(a)=f(b)$可知最大值和最小值中至少有一个是被$f$在内点$\xi\in(a,b)$处所取得.由$Fermat$定理,$f'(\xi)=0$.

\begin{ti}
\textbf{定理5.9($Lagrange$中值定理)}  设函数$f$在$[a,b]$上连续,在$(a,b)$内可导,则存在$\xi\in(a,b)$,使得
$$f(b)-f(a)=f'(\xi)(b-a)$$.
\end{ti}

\textbf{证明}:记$\nu=\dfrac{f(b)-f(a)}{b-a}$,$g(x)=f(x)-\nu x$,则$g(b)=g(a)$.根据$Rolle$定理,存在$\xi\in(a,b)$,使得$g'(\xi)=0$.此时$f'(\xi)=\nu$,即$f(b)-f(a)=f'(\xi)(b-a)$.

\begin{ti}
\textbf{定理5.10($Cauchy$中值定理)}  设函数$f,g$在$[a,b]$上连续,在$(a,b)$内可导,且$g'(x)\neq 0$,则存在$\xi\in(a,b)$,使得
$$\frac{f(b)-f(a)}{g(b)-g(a)}=\frac{f'(\xi)}{g'(\xi)}$$.
\end{ti}

\begin{ti}
\textbf{定义5.11(凸函数)}  设函数$f$定义在区间$I$上.如果对任意的$x,y\in I$和任意的$t\in[0,1]$,均有
$$f((1-t)x+ty)\leq(1-t)f(x)+tf(y)$$
成立,则称$f$为$I$上的\textbf{凸函数};若当$x\neq y$且$t\in(0,1)$时不等号严格成立,则称$f$为\textbf{严格凸函数}.
\end{ti}

\begin{ti}
\textbf{定理5.12($Jensen$不等式)}  设函数$f$为区间$I$上的凸函数.对任意的$x_1,\ldots,x_n\in I$以及满足$\displaystyle\sum_{i=1}^{n}\lambda_i=1$的非负实数$\lambda_i$,有
$$f\left(\sum_{i=1}^{n}\lambda_i x_i\right)\leq\sum_{i=1}^{n}\lambda_i f(x_i)$$.
\end{ti}

\begin{ti}
\textbf{定理5.13($\text{L'Hospital}$法则)} \\
(1)(之一)设$f,g$在$(a,b)$中可导,且$g'$处处非零.如果$\displaystyle\lim_{x\to a^+}f(x)=\lim_{x\to a^+}g(x)=0$,且$\displaystyle\lim_{x\to a^+}\dfrac{f'(x)}{g'(x)}=\alpha$,其中$\alpha$为实数或$\pm\infty$,则$\displaystyle\lim_{x\to a^+}\dfrac{f(x)}{g(x)}=\alpha$.\\
(2)(之二)设$f,g$在$(a,b)$中可导,且$g'$处处非零.如果$\displaystyle\lim_{x\to a^+}g(x)=\infty$,且$\displaystyle\lim_{x\to a^+}\dfrac{f'(x)}{g'(x)}=\alpha$,其中$\alpha$为实数或$\pm\infty$,则$\displaystyle\lim_{x\to a^+}\dfrac{f(x)}{g(x)}=\alpha$.
\end{ti}

\textbf{证明}:(1)补充定义$f(a)=g(a)=0$,则$f,g\in C^0[a,b)$.由$Cauchy$中值定理,任给$x\in(a,b)$,存在$\xi\in(a,x)$,使得
$$\frac{f(x)}{g(x)}=\frac{f(x)-f(a)}{g(x)-g(a)}=\frac{f'(\xi)}{g'(\xi)}.$$
当$x\to a^+$时,$\xi\to a^+$,从而$\displaystyle\lim_{x\to a^+}\frac{f(x)}{g(x)}=\lim_{x\to a^+}\frac{f'(\xi)}{g'(\xi)}=\alpha$. \\
(2)分情况讨论.\\
\romannum{1}  当$\alpha=0$时.
根据题设,任给$\varepsilon>0$,存在$\eta>0$,当$x\in(a,a+\eta)$时
$$\left\lvert\dfrac{f'(x)}{g'(x)}\right\rvert<\dfrac{\varepsilon}{2}.$$
此时,根据$Cauchy$中值定理,存在$\xi\in(x,a+\eta)$,使得
$$\dfrac{f(x)-f(a+\eta)}{g(x)-g(a+\eta)}=\dfrac{f'(\xi)}{g'(\xi)}$$.上式可以改写为
$$\frac{f(x)}{g(x)}=\frac{f'(\xi)}{g'(\xi)}\left[1-\frac{g(a+\eta)}{g(x)}\right]+\frac{f(a+\eta)}{g(x)}.$$
利用$\left\lvert\dfrac{f'(x)}{g'(x)}\right\rvert<\dfrac{\varepsilon}{2}$以及$\displaystyle\lim_{x\to a^+}g(x)=\infty$不难得知,
存在正数$\delta<\eta$,使得当$x\in(a,a+\delta)$时
$$\left\lvert\dfrac{f(x)}{g(x)}\right\rvert<\varepsilon,$$这说明$\displaystyle\lim_{x\to a^+}\dfrac{f(x)}{g(x)}=0$.\\
\romannum{2}  当$\alpha\in\mathbb{R}$时,通过将$f$换成$f-\alpha g$可以转化为$\alpha=0$的情形;\\
\romannum{3}  当$\alpha=\pm\infty$时与$\alpha=0$情形的证明类似.

\begin{ti}
\textbf{定理5.14($Taylor$公式,带$Peano$余项)}  设函数$f$在$x_0$处有$n$阶导数,则当$x\to x_0$时
$$f(x)=\sum_{k=0}^{n}\frac{f^{(k)}(x_0)}{k!}(x-x_0)^{k}+o\bigl((x-x_0)^{n}\bigr),$$
其中我们约定$f^{(0)}=f$,$(x-x_0)^{0}=1$.
\end{ti}

\begin{ti}
\textbf{定理5.15($Taylor$公式,带$Lagrange$余项)}  设$n\geq1$,$f$在包含$x_0$与$x$的某个开区间内$n$阶可导,则
$$f(x)=\sum_{k=0}^{n-1}\frac{f^{(k)}(x_0)}{k!}(x-x_0)^{k}+\frac{f^{(n)}(\xi)}{n!}(x-x_0)^{n},$$
其中$\xi=x_0+\theta(x-x_0)$,$\theta\in(0,1)$.
\end{ti}

\textbf{证明}:当$x=x_0$时结论显然成立.当$x\neq x_0$时,令$F(t)=\displaystyle\sum_{k=0}^{n-1}\frac{f^{(k)}(t)}{k!}(x-t)^k$,$G(t)=(x-t)^n$.逐项求导并将相邻项抵消,得到
$$F'(t)=\frac{f^{(n)}(t)}{(n-1)!}(x-t)^{n-1},\qquad G'(t)=-n(x-t)^{n-1}.$$
由$Cauchy$微分中值定理,存在$x_0$与$x$之间的$\xi$,使得
$$\frac{F(x)-F(x_0)}{G(x)-G(x_0)}=\frac{F'(\xi)}{G'(\xi)}=-\frac{f^{(n)}(\xi)}{n!}.$$
注意$F(x)=f(x)$,$G(x)=0$,$G(x_0)=(x-x_0)^n$,整理即得结论.

\subsection{一元函数的积分}
\begin{ti}
\textbf{定理5.16($Newton\text{-}Leibniz$公式)}  设函数$f$在$[a,b]$上连续,且函数$F$为$f$的一个原函数,则
$$\int_a^b f(x)\,\mathrm{d}x=F(b)-F(a)$$.
\end{ti}

\textbf{证明}:记$h(x)=\displaystyle\int_a^x f(t)\,\mathrm{d}t$,由变上限积分定理可知$h'=f$.于是$(F-h)'=0$处处成立,故$F-h$为常值函数.特别地,有
$$\int_a^b f(x)\,\mathrm{d}x=h(b)-h(a)=F(b)-F(a).$$

\begin{ti}
\textbf{命题5.17(分部积分法)}  设函数$u,v$在$[a,b]$上有连续导数,则
$$\int_a^b u(x)v'(x)\,\mathrm{d}x=u(x)v(x)\Big|_a^b-\int_a^b u'(x)v(x)\,\mathrm{d}x$$.
\end{ti}

\begin{ti}
\textbf{命题5.18(换元积分法)}  设$\varphi\in C^1[a,b]$,函数$f$在区间$I\supset\varphi([a,b])$中连续,则
$$\int_a^b f(\varphi(x))\varphi'(x)\,\mathrm{d}x=\int_{\varphi(a)}^{\varphi(b)}f(y)\,\mathrm{d}y$$.
\end{ti}

\begin{ti}
\textbf{定理5.19($Wallis$公式)} $\displaystyle\lim_{n\to\infty}\dfrac{1}{2n+1}\left[\dfrac{(2n)!!}{(2n-1)!!}\right]^{2}=\dfrac{\pi}{2}.$
\end{ti}
\textbf{证明}:当$m$为非负整数时,记$\displaystyle I_m=\int_0^{\frac{\pi}{2}}\sin^{m}x\,\mathrm{d}x$,则$I_0=\dfrac{\pi}{2}$,$I_1=1$.当$m\geq 2$时,利用分部积分可得
$$I_m=\int_0^{\frac{\pi}{2}}\sin^{m}x\,\mathrm{d}x=(m-1)\int_0^{\frac{\pi}{2}}\sin^{m-2}x\cos^{2}x\,\mathrm{d}x=(m-1)(I_{m-2}-I_m),$$
这就得到了递推公式$I_m=\dfrac{m-1}{m}I_{m-2}$,从而有
$$I_{2n}=\frac{2n-1}{2n}I_{2n-2}=\cdots=\frac{(2n-1)!!}{(2n)!!}\cdot\frac{\pi}{2},$$
同理可得
$$I_{2n+1}=\frac{(2n)!!}{(2n+1)!!}.$$
注意到当$0<x<\dfrac{\pi}{2}$时,$\sin^{2n+1}x<\sin^{2n}x<\sin^{2n-1}x$,由积分的保序性质可得$I_{2n+1}<I_{2n}<I_{2n-1}$,即
$$\frac{(2n)!!}{(2n+1)!!}<\frac{(2n-1)!!}{(2n)!!}\cdot\frac{\pi}{2}<\frac{(2n-2)!!}{(2n-1)!!}.$$
整理以后可得
$$\frac{n\pi}{2n+1}<\frac{1}{2n+1}\left[\frac{(2n)!!}{(2n-1)!!}\right]^{2}<\frac{\pi}{2}.$$
由夹逼定理可得
$$\lim_{n\to\infty}\frac{1}{2n+1}\left[\frac{(2n)!!}{(2n-1)!!}\right]^{2}=\frac{\pi}{2}.$$

\begin{ti}
\textbf{定理5.20($Stirling$公式)} $n!=\sqrt{2\pi n}\left(\dfrac{n}{e}\right)^n e^{\delta_n},\;0<\delta_n<\dfrac{1}{12n}$.
\end{ti}
\textbf{证明}:首先易得$\dfrac{1}{k+1}<\ln(1+\dfrac{1}{k})<\dfrac{1}{k},\;\forall k\geq 1$,上式可以改写为
$$\dfrac{1}{k+1}<\int_{k}^{k+1} \dfrac{1}{x}\,dx<\dfrac{1}{k},$$
考虑区间$[k,k+1]$的中点,记为$\alpha_k=k+\dfrac{1}{2}$.先后作变量替换$x=\alpha_k+t,\,t=\alpha_k s$可得
$$\int_{k}^{k+1} \dfrac{1}{x}\,dx=\int_{-\frac{1}{2}}^{\frac{1}{2}} \dfrac{1}{\alpha_k+t}\,dt=\int_{-\frac{1}{2\alpha_k}}^{\frac{1}{2\alpha_k}}\dfrac{1}{1+s}\,ds ,$$
即
$$\ln(1+\dfrac{1}{k})=\ln(1+\dfrac{1}{2\alpha_k})-\ln(1-\dfrac{1}{2\alpha_k}).$$
在$\ln(1\pm x)$的$Taylor$展开中代入$(2\alpha_k)^{-1}=(2k+1)^{-1}$可得
\begin{equation}
\ln(1+\dfrac{1}{k})=2\sum_{m=0}^{\infty}\dfrac{(2k+1)^{-2m-1}}{2m+1}.
\end{equation}
当$k\geq 1$时,记
$$c_k=\alpha_k\ln(1+\dfrac{1}{k})-1=(k+\frac{1}{2})\ln(1+\dfrac{1}{k})-1,$$
由(1)式可得
$$0<c_k=\sum_{m=1}^{\infty}\dfrac{(2k+1)^{-2m}}{2m+1}<\dfrac{1}{3}\sum_{m=1}^{\infty}(2k+1)^{-2m}=\dfrac{1}{12k(k+1)}.$$
记$\displaystyle C=\sum_{k=1}^{\infty}c_k=\lim_{k\to\infty}\sum_{l=1}^{k}c_l,$
由上式可知$0<C<\dfrac{1}{12}$.同理,记$\displaystyle\delta_n=C-\sum_{k=1}^{n-1}c_k$,则
$$0<\delta_n=\sum_{k=n}^{\infty}c_k<\dfrac{1}{12}\sum_{k=n}^{\infty}(\dfrac{1}{k}-\dfrac{1}{k+1})=\dfrac{1}{12n}.$$
另一方面,注意到
$$e^{1+c_k}=\left(1+\dfrac{1}{k}\right)^{k+\frac{1}{2}}=\dfrac{(k+1)^{k+1}}{k^k}\dfrac{\sqrt{k+1}}{\sqrt{k}}\dfrac{1}{k+1},$$
这说明
$$e^{C-\delta_n}=\prod_{k=1}^{n-1}e^{c_k}=\prod_{k=1}^{n-1}\left[(1+\dfrac{1}{k})^{k+\frac{1}{2}}e^{-1}\right]=n^n\sqrt{n}\dfrac{1}{n!}e^{1-n},$$
上式可以改写为
$$n!=e^{1-C}\sqrt{n}\left(\dfrac{n}{e}\right)^n e^{\delta_n}.$$
下面我们来确定$e^{1-C}$的值.首先利用上式可得
$$\lim_{n\to\infty}\dfrac{1}{\sqrt{2n+1}}\dfrac{(2n)!!}{(2n-1)!!}=\lim_{n\to\infty}\dfrac{1}{\sqrt{2n+1}}\dfrac{[2^n\,n!]^2}{(2n)!}=\dfrac{e^{1-C}}{2},$$
代入\textbf{定理5.19}的$Wallis$公式可得$e^{1-C}=\sqrt{2\pi}$.最后我们得到如下$Stirling$公式
\begin{equation}
n!=\sqrt{2\pi n}\left(\dfrac{n}{e}\right)^n e^{\delta_n},\;0<\delta_n<\dfrac{1}{12n}.
\end{equation}

\subsection{\texorpdfstring{$Riemann$}{Riemann}积分}
\begin{ti}
\textbf{定义5.21($Riemann$可积)}  设$f$为$[a,b]$中的有界函数.如果存在实数$I$,使得任给$\varepsilon>0$,均存在$\delta>0$,当分割$P$的模满足$\left\lVert P\right\rVert<\delta$时,均有
$$\left\lvert\sum_{i=1}^{n}f(\xi_i)\Delta x_i-I\right\rvert<\varepsilon,\quad \forall\,\xi_i\in[x_{i-1},x_i],$$
则称$f$在$[a,b]$中\textbf{$Riemann$可积}(简称可积),记为$f\in R[a,b]$. $I$称为$f$在$[a,b]$中的\textbf{$Riemann$积分}(简称积分),记为
$$\int_a^b f(x)\,\mathrm{d}x=I=\lim_{\left\lVert P\right\rVert\to 0}\sum_{i=1}^{n}f(\xi_i)\Delta x_i.$$
\end{ti}

\begin{ti}
\textbf{定理5.22(可积的充要条件)}  设函数$f$在$[a,b]$上有界,则以下命题等价: \\
(1) $f$在$[a,b]$中$Riemann$可积; \\
(2) $f$在$[a,b]$中上积分和下积分相等; \\
(3) $\displaystyle\lim_{\left\lVert P\right\rVert\to 0}\sum_{i=1}^{n}\omega_i\varDelta x_i=0$; \\
(4) 任给$\varepsilon>0,$存在$[a,b]$中某个分割$P$,使得
$$S(P)-s(P)=\sum_{i=1}^{n}\omega_i\varDelta x_i<\varepsilon$$
\end{ti}

\begin{ti}
\textbf{定义5.23(零测集)}  设$A\subset\mathbb{R}$.如果任给$\varepsilon>0$,均可找到至多可数个开区间$\{I_i\}$,使得$A$包含于这些区间之并,且
$$\sum_{i\leq n}\left\lvert I_i\right\rvert\leq\varepsilon,\quad \forall\,n\geq 1,$$
即这些区间的长度之和不超过$\varepsilon$,则称$A$为\textbf{零测集}.
\end{ti}

\begin{ti}
\textbf{定理5.24($Lebesgue$可积性判别法)}  设函数$f$在$[a,b]$上有界,则$f$在$[a,b]$上$Riemann$可积当且仅当$f$的间断点集为零测集.
\end{ti}

\begin{ti}
\textbf{命题5.25}  设$f\in C^0[a,b]$,$c\in[a,b]$.定义
$$g(x)=\int_c^x f(t)\,\mathrm{d}t,\quad x\in[a,b],$$
则$g'=f$在$[a,b]$中处处成立.
\end{ti}

\begin{ti}
\textbf{命题5.26(阶梯逼近)}  设$f\in R[a,b]$,则存在一列阶梯函数$f_n$,使得
$$\lim_{n\to\infty}\int_a^b|f_n(x)-f(x)|\,\mathrm{d}x=0,$$
且每一个$f_n$均介于$f$的上、下确界之间.此外,任给$g\in R[a,b]$,均有
$$\lim_{n\to\infty}\int_a^b f_n(x)g(x)\,\mathrm{d}x=\int_a^b f(x)g(x)\,\mathrm{d}x.$$
\end{ti}

\begin{ti}
\textbf{命题5.27(分段线性逼近)}  设$f\in R[a,b]$,则存在一列连续的分段线性函数$f_n$,使得$f_n(a)=f(a)$,$f_n(b)=f(b)$,且
$$\lim_{n\to\infty}\int_a^b|f_n(x)-f(x)|\,\mathrm{d}x=0,$$
并且每一个$f_n$均介于$f$的上、下确界之间.此外,任给$g\in R[a,b]$,均有
$$\lim_{n\to\infty}\int_a^b f_n(x)g(x)\,\mathrm{d}x=\int_a^b f(x)g(x)\,\mathrm{d}x.$$
\end{ti}

\begin{ti}
\textbf{命题5.28($Riemann\text{-}Lebesgue$引理)}  设$f\in R[a,b]$,则
$$\lim_{\lambda\to+\infty}\int_a^b f(x)\sin(\lambda x)\,\mathrm{d}x=0,\qquad \lim_{\lambda\to+\infty}\int_a^b f(x)\cos(\lambda x)\,\mathrm{d}x=0.$$
\end{ti}

\begin{ti}
\textbf{定理5.29(分部积分之二)}  设$f,g\in R[a,b]$,$C_1,C_2\in\mathbb{R}$.当$x\in[a,b]$时,记
$$F(x)=\int_a^x f(t)\,\mathrm{d}t+C_1,\quad G(x)=\int_a^x g(t)\,\mathrm{d}t+C_2,$$
则有
$$\int_a^b F(x)g(x)\,\mathrm{d}x=F(x)G(x)\Big|_a^b-\int_a^b f(x)G(x)\,\mathrm{d}x.$$
\end{ti}

\textbf{证明}:如果$f,g\in C^0[a,b]$,则$F,G\in C^1[a,b]$,欲证等式可从已知的分部积分公式直接得到.一般地,我们分别取$f,g$的分段线性逼近$\{f_n\},\{g_n\}$,使得
$$\int_a^b|f_n(x)-f(x)|\,\mathrm{d}x<\frac{1}{n},\quad \int_a^b|g_n(x)-g(x)|\,\mathrm{d}x<\frac{1}{n}.$$
当$x\in[a,b]$时,记
$$F_n(x)=\int_a^x f_n(t)\,\mathrm{d}t+C_1,\quad G_n(x)=\int_a^x g_n(t)\,\mathrm{d}t+C_2.$$
由$f_n,g_n$连续可得
$$\int_a^b F_n(x)g_n(x)\,\mathrm{d}x=F_n(x)G_n(x)\Big|_a^b-\int_a^b f_n(x)G_n(x)\,\mathrm{d}x.$$
当$n\to\infty$时,我们来研究上式中的各项如何变化.首先,不妨设$f,g,f_n,g_n,F_n,G_n$在$[a,b]$中的绝对值均不超过$M$.此时有
$$|F_n(x)G_n(x)-F(x)G(x)|\leq|F_n(x)[G_n(x)-G(x)]|+|[F_n(x)-F(x)]G(x)|$$
$$\leq M\int_a^x|g_n(t)-g(t)|\,\mathrm{d}t+M\int_a^x|f_n(t)-f(t)|\,\mathrm{d}t<\frac{2M}{n}.$$
上式中其余的项可以类似地作估计.令$n\to\infty$就由上面的等式得到了欲证等式.

\begin{ti}
\textbf{定理5.30(积分第二中值定理)}  设函数$f$在$[a,b]$上可积,$g$在$[a,b]$上单调,则存在$\xi\in[a,b]$,使得
$$\int_a^b f(x)g(x)\,\mathrm{d}x=g(a)\int_a^{\xi}f(x)\,\mathrm{d}x+g(b)\int_{\xi}^{b}f(x)\,\mathrm{d}x$$.
\end{ti}

\textbf{证明}:记$F(x)=\displaystyle\int_a^x f(t)\,\mathrm{d}t$,则$F(a)=0$且$F$在$[a,b]$上连续.若$g$为常数,结论显然.先设$g$单调递减.对分割$P:a=x_0<\cdots<x_m=b$,令$g_P(t)=g(x_{i-1})$($t\in(x_{i-1},x_i]$),并记
$$A_P=\int_a^b f(t)g_P(t)\,\mathrm{d}t=\sum_{i=1}^m g(x_{i-1})[F(x_i)-F(x_{i-1})].$$
由有限求和分部公式,
$$A_P=g(b)F(b)+\sum_{i=1}^m[g(x_{i-1})-g(x_i)]F(x_i).$$
后一个和式中的权重非负,总和为$g(a)-g(b)$.由$F$的连续性及介值定理,存在$\xi_P\in[a,b]$,使它等于$[g(a)-g(b)]F(\xi_P)$.又由单调性,
$$\int_a^b|g_P(t)-g(t)|\,\mathrm{d}t\leq\lVert P\rVert\,[g(a)-g(b)],$$
故当$\lVert P\rVert\to0$时,$A_P\to\int_a^b fg$.从$\{\xi_P\}$中取收敛子列,再用$F$的连续性,得$\int_a^bfg=g(b)F(b)+[g(a)-g(b)]F(\xi)$,即所求公式.若$g$单调递增,对$-g$应用上述结论即可.

\begin{ti}
\textbf{推论5.31($Bonnet$)}  设$f\in R[a,b]$,$g$为非负函数.\\
(1)如果$g$在$[a,b]$中单调递减,则存在$\zeta\in[a,b]$,使得
$$\int_a^b f(x)g(x)\,\mathrm{d}x=g(a)\int_a^{\zeta}f(x)\,\mathrm{d}x;$$
(2)如果$g$在$[a,b]$中单调递增,则存在$\eta\in[a,b]$,使得
$$\int_a^b f(x)g(x)\,\mathrm{d}x=g(b)\int_{\eta}^{b}f(x)\,\mathrm{d}x.$$
\end{ti}

\begin{ti}
\textbf{命题5.32($Euler\text{-}Poisson$积分)} $\;\displaystyle\int_{0}^{+\infty} e^{-x^2}\,dx=\dfrac{\sqrt{\pi}}{2}$.
\end{ti}

\textbf{解法一}:先计算积分$\displaystyle I_n=\int_{0}^{+\infty} \dfrac{1}{(1+x^2)^n}\,dx $. \\
当$n=1$时,$\displaystyle I_1=\left.\arctan x\right|_0^{+\infty}=\dfrac{\pi}{2}.$当$n\geq 2$时,
$$I_n=\int_{0}^{+\infty} \dfrac{1+x^2-x^2}{(1+x^2)^n}\,dx=I_{n-1}-\int_{0}^{+\infty} \dfrac{x^2}{(1+x^2)^n}\,dx $$
$$\phantom{I_n=}=I_{n-1}+\dfrac{1}{2(n-1)}\left.\dfrac{x}{(1+x^2)^{n-1}}\right|_0^{+\infty}-\dfrac{1}{2(n-1)}\int_{0}^{+\infty} \dfrac{1}{(1+x^2)^{n-1}}\,dx $$
$$=\dfrac{2n-3}{2n-2}I_{n-1}.$$
这说明$I_n=\dfrac{\pi}{2}\dfrac{(2n-3)!!}{(2n-2)!!}$. \\

利用不等式$e^t\geq 1+t$可得$1-x^2\leq e^{-x^2}\leq(1+x^2)^{-1}$.当$n\geq 1$时就有
$$\int_{0}^{1} (1-x^2)^n\,dx\leq\int_{0}^{1} e^{-nx^2}\,dx\leq\int_{0}^{+\infty} e^{-nx^2}\,dx\leq\int_{0}^{+\infty} \dfrac{1}{(1+x^2)^n}\,dx,$$
上式左边利用变量代换$x=\cos t$以及分部积分,中间利用$x=\dfrac{t}{\sqrt{n}}$,右边利用上面计算得到的积分,得到
$$\dfrac{(2n)!!}{(2n+1)!!}\leq\dfrac{1}{\sqrt{n}}\int_{0}^{+\infty} e^{-t^2}\,dt\leq\dfrac{\pi}{2}\dfrac{(2n-3)!!}{(2n-2)!!},$$
两边乘$\sqrt{n}$,再利用$Wallis$公式不难看出,当$n\to\infty$时,左右极限均为$\dfrac{\sqrt{\pi}}{2}$,这说明
$$\int_{0}^{+\infty} e^{-t^2}\,dt=\dfrac{\sqrt{\pi}}{2} .$$

\textbf{解法二}:利用重积分与极坐标变换.\\
$$I^2=\left(\int_{0}^{+\infty}  e^{-x^2}\,dx \right)\left(\int_{0}^{+\infty}  e^{-y^2}\,dy \right)=\iint\limits_{\substack{x\in[0,+\infty) \\ y\in[0,+\infty)}} e^{-(x^2+y^2)}\,dxdy$$
$$=\iint\limits_{\substack{r\in[0,+\infty) \\ \theta\in[0,\frac{\pi}{2}]}}re^{-r^2}\,drd\theta=\dfrac{\pi}{2}\int_{0}^{+\infty} re^{-r^2}\,dr=\dfrac{\pi}{4}.$$
于是有
$$I=\dfrac{\sqrt{\pi}}{2}.$$

\begin{ti}
\textbf{命题5.33}  计算积分 $\displaystyle\int_{0}^{+\infty}\dfrac{\sin^2 x}{x^2}\,\mathrm{d}x$.
\end{ti}

\textbf{证明}:当 $n\geq 0$ 时,利用等式 $\sin(2n+1)x-\sin(2n-1)x=2\sin x\cos 2nx$ 易得
\[\int_{0}^{\frac{\pi}{2}}\dfrac{\sin(2n+1)x}{\sin x}\,\mathrm{d}x=\dfrac{\pi}{2}.\tag{1}\]
再利用等式 $\sin^2 nx-\sin^2(n-1)x=\sin(2n-1)x\sin x$ 可得
\[\int_{0}^{\frac{\pi}{2}}\dfrac{\sin^2 nx}{\sin^2 x}\,\mathrm{d}x=\sum_{k=1}^{n}\int_{0}^{\frac{\pi}{2}}\dfrac{\sin(2k-1)x}{\sin x}\,\mathrm{d}x=\dfrac{\pi}{2}n.\tag{2}\]
当 $x\in\left(0,\dfrac{\pi}{2}\right)$ 时,注意到
\[0<\dfrac{1}{\sin^2 x}-\dfrac{1}{x^2}=1+\dfrac{1}{\tan^2 x}-\dfrac{1}{x^2}<1,\tag{3}\]
因此有
\begin{align*}
\int_0^{+\infty}\dfrac{\sin^2 x}{x^2}\,\mathrm{d}x
&=\lim_{n\to\infty}\dfrac{1}{n}\int_0^{+\infty}\dfrac{\sin^2 nx}{x^2}\,\mathrm{d}x=\lim_{n\to\infty}\dfrac{1}{n}\int_0^{\frac{\pi}{2}}\dfrac{\sin^2 nx}{x^2}\,\mathrm{d}x\\
&=\lim_{n\to\infty}\dfrac{1}{n}\int_0^{\frac{\pi}{2}}\sin^2 nx\left(\dfrac{1}{x^2}-\dfrac{1}{\sin^2 x}\right)\mathrm{d}x+\lim_{n\to\infty}\dfrac{1}{n}\int_0^{\frac{\pi}{2}}\dfrac{\sin^2 nx}{\sin^2 x}\,\mathrm{d}x\\
&=0+\dfrac{\pi}{2}=\dfrac{\pi}{2}.
\end{align*}
\textbf{注}:\;注意到 $x^{-2}=(-x^{-1})'$, 利用分部积分和简单的变量替换,由此例不难得出
\[\int_{0}^{+\infty}\dfrac{\sin^2 x}{x^2}\,\mathrm{d}x=\int_{0}^{+\infty}\dfrac{\sin x}{x}\,\mathrm{d}x=\dfrac{\pi}{2}.\tag{4}\]

\section{多元微分学与多元积分学}
\subsection{多元函数的微分}
\begin{ti}
\textbf{定义6.1(方向导数)}  设函数$f$在点$\mathbf{x}_0\in\mathbb{R}^n$的某邻域中有定义,$\mathbf{v}$为单位向量.如果极限
$$\lim_{t\to 0}\frac{f(\mathbf{x}_0+t\mathbf{v})-f(\mathbf{x}_0)}{t}$$
存在,则称其为$f$在$\mathbf{x}_0$处沿方向$\mathbf{v}$的\textbf{方向导数},记为$\dfrac{\partial f}{\partial\mathbf{v}}(\mathbf{x}_0)$.
\end{ti}

\begin{ti}
\textbf{定义6.2(微分)}  设$D\subset \mathbb{R}^n$为开集,$f:D\rightarrow \mathbb{R}$为多元函数,
$x^0\in D$.如果存在线性映射$L:\mathbb{R}^n\rightarrow\mathbb{R}$,使得在$x^0$附近成立
$$f(x)-f(x^0)=L(x-x^0)+o(\left\lVert x-x^0\right\rVert)\;(x\rightarrow x^0),$$
成立,则称$f$在$x^0$处\textbf{可微},称$L$为$f$在$x^0$处的\textbf{微分},也记为$df(x^0)$
\end{ti}

\begin{ti}
\textbf{命题6.3}  设函数$f$在$\mathbf{x}_0$处可微,则$f$在$\mathbf{x}_0$处连续,且$f$在$\mathbf{x}_0$处的方向导数都存在,并且有
$$\frac{\partial f}{\partial \mathbf{u}}(\mathbf{x}_0)=df(\mathbf{x}_0)(\mathbf{u}),$$
其中$\mathbf{u}$为单位向量.
\end{ti}

\begin{ti}
\textbf{命题6.4(可微的充分条件)}  设函数$f$在$\mathbf{x}_0$的某邻域内各偏导数均存在且在$\mathbf{x}_0$处连续,则$f$在$\mathbf{x}_0$处可微.
\end{ti}

\begin{ti}
\textbf{定义6.5(梯度)}  设函数$f$在$\mathbf{x}_0$处可微,记
$$\nabla f(\mathbf{x}_0)=\left(\frac{\partial f}{\partial x_1}(\mathbf{x}_0),\ldots,\frac{\partial f}{\partial x_n}(\mathbf{x}_0)\right),$$
称为$f$在$\mathbf{x}_0$处的\textbf{梯度}.利用$\mathbb{R}^n$中的标准内积,有$df(\mathbf{x}_0)(\mathbf{v})=\nabla f(\mathbf{x}_0)\cdot\mathbf{v}$.
\end{ti}

\begin{ti}
\textbf{定理6.6(链式法则)}  设$D$和$\varDelta$分别为$\mathbb{R}^n\text{和}\mathbb{R}^m$
中的开集,$f:D\rightarrow \mathbb{R}^m\text{和}g:\varDelta\rightarrow\mathbb{R}^l$为向量值函数,且$f(D)\subset\varDelta.$
如果$f$在$x^0\in D$处可微,$g$在$y^0=f(x^0)$处可微,则复合函数$h=g\circ f$在$x^0$处可微,且
$$Jh(x^0)=Jg(y^0)\cdot Jf(x^0)$$.
\end{ti}

\begin{ti}
\textbf{定理6.7(微分中值定理)} 设$D\subset \mathbb{R}^n$为凸域,$f:D\rightarrow\mathbb{R}$在$D$中处处可微,则任给$x,y\in D$,
存在$\theta\in(0,1)$,使得
$$f(x)-f(y)=\nabla f(\xi)\cdot(x-y),\;\xi=\theta x+(1-\theta)y$$.
\end{ti}

\begin{ti}
\textbf{定理6.8(拟微分中值定理)}  设$D\subset\mathbb{R}^n$为凸域,$f:D\rightarrow\mathbb{R}^m$在$D$中处处可微,则对任意的$x,y\in D$,存在$\xi\in D$,使得
$$\left\lVert f(x)-f(y)\right\rVert\leq\left\lVert Jf(\xi)\right\rVert\cdot\left\lVert x-y\right\rVert.$$
\end{ti}

\textbf{证明}:基本的想法是对$f$的分量的线性组合应用微分中值定理.为此,不妨设$f(x)\neq f(y)$.
任意取定$\mathbb{R}^m$中的单位向量$u=(u_1,\ldots,u_m)$,记
$$g=u\cdot f=\displaystyle\sum_{i=1}^{m}u_if_i,$$
则$g$为$D$中可微函数.根据微分中值定理,存在$\xi\in D$,使得
$$g(x)-g(y)=\nabla g(\xi)\cdot(x-y).$$
注意到$\nabla g(\xi)=\displaystyle\sum_{i=1}^{m}u_i\nabla f_i(\xi)$,利用$Cauchy-Schwarz$不等式可得
$$\left\lVert\nabla g(\xi)\right\rVert\leq\left\lVert Jf(\xi)\right\rVert.$$
由$g(x)-g(y)=u\cdot[f(x)-f(y)]$可得
$$\left\lvert u\cdot[f(x)-f(y)]\right\rvert\leq\left\lVert Jf(\xi)\right\rVert\cdot\left\lVert x-y\right\rVert.$$
在上式中取$u=\dfrac{f(x)-f(y)}{\left\lVert f(x)-f(y)\right\rVert}$就完成了定理的证明.

\begin{ti}
\textbf{定理6.9(逆映射定理)}  设$D\subset\mathbb{R}^n$为开集,$f:D\rightarrow\mathbb{R}^n$为$C^k\,(k\geq 1)$映射,$\mathbf{x}_0\in D$.
如果$\det Jf(\mathbf{x}_0)\neq 0$,则存在$\mathbf{x}_0$的开邻域$U\subset D$以及$y_0=f(\mathbf{x}_0)$的开邻域$V\subset\mathbb{R}^n$,使得$f|_U:U\rightarrow V$是可逆映射,且其逆仍为$C^k$映射.
\end{ti}

\textbf{证明}:不失一般性,可设$\mathbf{x}_0=0$,$y_0=0$.以$L$记$f$在$\mathbf{x}_0=0$处的微分,则$L$可逆,且$L^{-1}\circ f$在$\mathbf{x}_0$处的微分为恒同映射.
如果欲证结论对$L^{-1}f$成立,则对$f$也成立.因此,不妨从一开始就假设$Jf(\mathbf{x}_0)=I_n$.\\

在$\mathbf{x}_0=0$附近,$f$是恒同映射的小扰动:$f(x)=x+g(x)$,$Jg(0)=0$,且$g(0)=0$.取$\delta>0$,使闭球$\overline B_\delta(0)\subset D$,并使$\left\lVert Jg(x)\right\rVert\leq\dfrac{1}{2}$在该闭球上成立.由\textbf{定理6.8}可知
$$\left\lVert g(x_1)-g(x_2)\right\rVert\leq\dfrac{1}{2}\left\lVert x_1-x_2\right\rVert,\quad x_1,x_2\in\overline B_\delta(0).$$
给定$y\in V=B_{\delta/2}(0)$,解$f(x)=y$等价于解$x=y-g(x)$.记$\varphi_y(x)=y-g(x)$,则当$x\in\overline B_\delta(0)$时,
$$\lVert\varphi_y(x)\rVert\leq\lVert y\rVert+\tfrac12\lVert x\rVert<\delta.$$
故$\varphi_y$是完备空间$\overline B_\delta(0)$到自身的压缩映射.由压缩映射原理,它有唯一不动点$x_y\in B_\delta(0)$.令$U=f^{-1}(V)\cap B_\delta(0)$,则$U$为开邻域,$f|_U:U\to V$为一一且满的映射,其逆映射$h(y)=x_y$满足$y-g(h(y))=h(y)$. \\

(1)$h:V\to U$是连续映射:当$y_1,y_2\in V$时,
$$\left\lVert h(y_1)-h(y_2)\right\rVert\leq\left\lVert y_1-y_2\right\rVert+\frac{1}{2}\left\lVert h(y_1)-h(y_2)\right\rVert,$$
这说明$\left\lVert h(y_1)-h(y_2)\right\rVert\leq 2\left\lVert y_1-y_2\right\rVert$. \\
(2)$h:V\to U$是可微映射:设$y_0\in V$,则对$y\in V$,有
$$h(y)-h(y_0)=(y-y_0)-[g(h(y))-g(h(y_0))]$$
$$\phantom{h(y)-h(y_0)}=(y-y_0)-Jg(h(y_0))\cdot(h(y)-h(y_0))+o(\left\lVert h(y)-h(y_0)\right\rVert).$$
利用$Jf=I_n+Jg$和(1),上式可改写为$Jf(h(y_0))\cdot(h(y)-h(y_0))=(y-y_0)+o(\left\lVert y-y_0\right\rVert)$,因而$h(y)-h(y_0)=[Jf(h(y_0))]^{-1}\cdot(y-y_0)+o(\left\lVert y-y_0\right\rVert)$.这说明$h$在$y_0$处可微. \\
(3)$h:V\to U$为$C^k$映射:由(2)可知$Jh(y)=[Jf(h(y))]^{-1}$,$\forall y\in V$.由$f\in C^k$知$Jf\in C^{k-1}$.由(2)及上式可推出$Jh$连续,即$h\in C^1$.再由$Jf\in C^{k-1}$,$h\in C^1$及上式可推出$Jh\in C^1$,即$h\in C^2$.依此类推,最后我们就得到$h\in C^k$.

\begin{ti}
\textbf{定理6.10(隐映射定理)}  设$W\subset \mathbb{R}^{m+n}$为开集,$W$中的点用$(x,y)$表示,
其中$x=(x_1,x_2,...,x_n),y=(y_1,y_2,...,y_m)$.设$k\geq1$,$f:W\rightarrow\mathbb{R}^m$为$C^k$映射,用分量表示为
$$f(x,y)=(f_1(x,y),f_2(x,y),...,f_m(x,y)).$$
设$(x^0,y^0)\in W,\text{且}det J_yf(x^0,y^0)\neq 0,\text{其中}J_yf(x,y)=\left(\dfrac{\partial f_i}{\partial y_j}(x,y)\right)_{m\times m}$.
则存在$x^0$的开邻域$V\subset\mathbb{R}^n$及$y^0$的开邻域$Y\subset\mathbb{R}^m$,使得对每个$x\in V$,方程$f(x,y)=f(x^0,y^0)$在$Y$中恰有一个解$y=\psi(x)$,且$\psi:V\to Y$为$C^k$映射.于是 \\
(1)$y^0=\psi(x^0),f(x,\psi(x))=f(x^0,y^0),\forall\,x\in V$; \\
(2)$J\psi(x)=-[J_yf(x,\psi(x))]^{-1}J_xf(x,\psi(x)),\text{其中}J_xf(x,y)=\left(\dfrac{\partial f_i}{\partial x_j}(x,y)\right)_{m\times n}$.
\end{ti}

\begin{ti}
\textbf{定义6.11(极值)}  设$D\subset\mathbb{R}^n$,$f:D\to\mathbb{R}$为多元函数,$\mathbf{x}_0\in D$.如果存在$\delta>0$,使得
$$f(\mathbf{x}_0)\geq f(\mathbf{x})\;\text{或}\;f(\mathbf{x}_0)\leq f(\mathbf{x}),\quad \forall\,\mathbf{x}\in D\cap B_\delta(\mathbf{x}_0)\setminus\{\mathbf{x}_0\},$$
则称$\mathbf{x}_0$为$f$的\textbf{极大(小)值点},$f(\mathbf{x}_0)$为$f$的\textbf{极大(小)值}.当上式中``$\geq$''(``$\leq$'')换成``$>$''(``$<$'')时,相应地把$\mathbf{x}_0$称为\textbf{严格极值点},$f(\mathbf{x}_0)$为\textbf{严格极值}.
\end{ti}

\begin{ti}
\textbf{命题6.12(极值的必要条件)}  设$\mathbf{x}_0$为函数$f$的极值点.如果$x_0$为$D$的内点,且$f$在$\mathbf{x}_0$处可微,则$\nabla f(\mathbf{x}_0)=\mathbf{0}$.
\end{ti}

\begin{ti}
\textbf{定理6.13(极值的充分条件)}  设函数$f$在$\mathbf{x}_0$的某邻域内二阶连续可微,且$\nabla f(\mathbf{x}_0)=\mathbf{0}$,
记$Hesse$矩阵$\nabla^2 f=\left(\dfrac{\partial^2 f}{\partial x_i\partial x_j}\right)_{n\times n}$.\\
(1)如果$\nabla^2 f(\mathbf{x}_0)$正定,则$\mathbf{x}_0$为严格极小值点;\\
(2)如果$\nabla^2 f(\mathbf{x}_0)$负定,则$\mathbf{x}_0$为严格极大值点;\\
(3)如果$\nabla^2 f(\mathbf{x}_0)$不定,则$\mathbf{x}_0$不是极值点.
\end{ti}

\begin{ti}
\textbf{定理6.14($Lagrange$乘数法)}  设$U\subset\mathbb{R}^n$为开集,$f:U\to\mathbb{R}$为可微函数,$\Phi:U\to\mathbb{R}^m\,(n>m)$为$C^k\,(k\geq 1)$映射.记$\Sigma=\Phi^{-1}(0)=\left\{\mathbf{x}\in U\;|\;\Phi(\mathbf{x})=0\right\}$.设$\mathbf{x}_0\in\Sigma$为$f$的条件极值点.如果$J\Phi(\mathbf{x}_0)$的秩为$m$,则存在$\lambda^0\in\mathbb{R}^m$,使得
$$\nabla f(\mathbf{x}_0)-\lambda^0\cdot J\Phi(\mathbf{x}_0)=0$$.
\end{ti}

\textbf{证明}:不妨设$\dfrac{\partial(\varphi_1,\varphi_2,\ldots,\varphi_m)}{\partial(x_1,x_2,\ldots,x_m)}(\mathbf{x}_0)\neq 0$,其中$\varphi_i$为$\Phi$的第$i$个分量.由隐映射定理,存在点$\mathbf{z}_0=(x_{m+1}^0,\ldots,x_n^0)$的开邻域$V$以及$C^k$映射$\psi:V\to\mathbb{R}^m$,使得$\mathbf{y}_0=\psi(\mathbf{z}_0)$,$\Phi(\psi(\mathbf{z}),\mathbf{z})=0$对一切$\mathbf{z}\in V$成立,其中$\mathbf{y}=(x_1,\ldots,x_m)$,$\mathbf{z}=(x_{m+1},\ldots,x_n)$,$\mathbf{x}=(\mathbf{y},\mathbf{z})\in U$.在$\mathbf{x}_0=(\mathbf{y}_0,\mathbf{z}_0)$处求导,有
$$J\psi(\mathbf{z}_0)=-(J_y\Phi)^{-1}(\mathbf{x}_0)\cdot J_z\Phi(\mathbf{x}_0).$$
由于$\mathbf{x}_0$为$f$的条件极值点,故$\mathbf{z}_0$为$f(\psi(\mathbf{z}),\mathbf{z})$的极值点(驻点),在$\mathbf{z}_0$处求导,得
$$J_yf(\mathbf{x}_0)\cdot J\psi(\mathbf{z}_0)+J_zf(\mathbf{x}_0)=0.$$
将$J\psi(\mathbf{z}_0)$的表达式代入上式,得
$$J_zf(\mathbf{x}_0)=J_yf(\mathbf{x}_0)\cdot(J_y\Phi)^{-1}(\mathbf{x}_0)\cdot J_z\Phi(\mathbf{x}_0).$$
记$\lambda^0=J_yf(\mathbf{x}_0)\cdot(J_y\Phi)^{-1}(\mathbf{x}_0)$,即$J_yf(\mathbf{x}_0)=\lambda^0\cdot J_y\Phi(\mathbf{x}_0)$,且上式可用$\lambda^0$改写为$J_zf(\mathbf{x}_0)=\lambda^0\cdot J_z\Phi(\mathbf{x}_0)$.两式结合起来就得到$\nabla f(\mathbf{x}_0)-\lambda^0\cdot J\Phi(\mathbf{x}_0)=0$.

\subsection{多元函数的积分}
\begin{ti}
\textbf{定义6.15(二重$Riemann$积分)}  设函数$f$在有界闭区域$D\subset\mathbb{R}^2$上有界.对$D$的任意分割$P=\{D_1,\ldots,D_n\}$,记
$$S(P)=\sum_{i=1}^{n}M_i\,\sigma(D_i),\qquad s(P)=\sum_{i=1}^{n}m_i\,\sigma(D_i)$$
其中$\displaystyle M_i=\sup_{D_i}f$,$\displaystyle m_i=\inf_{D_i}f$, $\sigma(D_i)$为$D_i$的面积.
如果$\displaystyle\inf_P S(P)=\sup_P s(P)$,则称$f$在$D$上\textbf{$Riemann$可积},
该公共值记为$\displaystyle\iint_D f(x,y)\,\mathrm{d}x\mathrm{d}y$.
\end{ti}

\begin{ti}
\textbf{命题6.16(重积分的交换积分次序)} 设$A\subset\mathbb{R}^2$为可求面积的有界集合,$f:A\rightarrow\mathbb{R}$为连续有界函数.
记$A$在$x$轴上的垂直投影为
$$I=\{x\in\mathbb{R}\;|\;\text{存在}y\text{使得}(x,y)\in A\}$$.
如果对于每一点$x\in I,A_x=\{y\in\mathbb{R}\;|\;(x,y)\in A\}$是区间(可退化为一点),则
$$\int_A f=\int_I dx\int_{A_x}f(x,y)dy.$$
\end{ti}

\begin{ti}
\textbf{命题6.17(重积分变量替换公式)}  设映射$\Phi:(u,v)\mapsto(x,y)$为有界闭区域$D'\to D$的$C^1$一一映射,且$\det D\Phi\neq 0$,函数$f$在$D$上连续,则
$$\iint_D f(x,y)\,\mathrm{d}x\mathrm{d}y=\iint_{D'}f(\Phi(u,v))\left\lvert\det D\Phi(u,v)\right\rvert\,\mathrm{d}u\mathrm{d}v.$$
特别地,极坐标变换$x=r\cos\theta, y=r\sin\theta$下,$\mathrm{d}x\mathrm{d}y=r\,\mathrm{d}r\mathrm{d}\theta$.
\end{ti}

\begin{ti}
\textbf{定义6.18(第一型曲线积分)}  设$\sigma:[\alpha,\beta]\rightarrow\mathbb{R}^n$为可求长的参数曲线,$s(t)$表示从$\sigma(\alpha)$到$\sigma(t)$的弧长,$f$是定义在曲线迹上的有界函数.
任给$[\alpha,\beta]$的分割$\pi:\alpha=t_0<t_1<\cdots<t_m=\beta$,取$\xi_i\in[t_{i-1},t_i]$,考虑和$\displaystyle\sum_{i=1}^{m}f(\sigma(\xi_i))\Delta s_i$,其中$\Delta s_i=s(t_i)-s(t_{i-1})$.如果极限
$$\lim_{\left\lVert\pi\right\rVert\to 0}\sum_{i=1}^{m}f(\sigma(\xi_i))\Delta s_i$$
存在且与$\{\xi_i\}$的选取无关,则称此极限为$f$在$\sigma$上的\textbf{第一型曲线积分},记为$\displaystyle\int_\sigma f\,\mathrm{d}s$.当$\sigma$为分段$C^1$曲线时,
$$\int_\sigma f\,\mathrm{d}s=\int_{\alpha}^{\beta}f(\sigma(t))\left\lVert\sigma'(t)\right\rVert\,\mathrm{d}t.$$
\end{ti}

\begin{ti}
\textbf{定义6.19(第二型曲线积分)} 对于分段$C^1$曲线,有
$$\int_{\sigma}f_1dx_1+f_2dx_2+...+f_ndx_n=\int_{\alpha}^{\beta}F(\sigma)\cdot\sigma'(t)\mathrm{d}t$$,
其中$F=(f_1,f_2,...,f_n),\sigma'(t)=(x_1'(t),x_2'(t),...,x_n'(t))$为曲线$\sigma$的切向量.
\end{ti}

\begin{ti}
\textbf{定义6.20(第一型曲面积分)}  设曲面$S$由参数方程$\mathbf{r}(u,v)$, $(u,v)\in D$给出,函数$f$在$S$上连续,则
$$\iint_S f\,\mathrm{d}S=\iint_D f(\mathbf{r}(u,v))\left\lVert\mathbf{r}_u\times\mathbf{r}_v\right\rVert\,\mathrm{d}u\mathrm{d}v$$.
\end{ti}

\begin{ti}
\textbf{定义6.21(第二型曲面积分)}  设向量场$\mathbf{F}=(P,Q,R)$在定向曲面$S$上连续,$\mathbf{n}$为$S$的单位法向量,则
$$\iint_S P\,\mathrm{d}y\wedge\mathrm{d}z+Q\,\mathrm{d}z\wedge\mathrm{d}x+R\,\mathrm{d}x\wedge\mathrm{d}y=\iint_S\mathbf{F}\cdot\mathbf{n}\,\mathrm{d}S$$.
\end{ti}

\begin{ti}
\textbf{定理6.22($Green$公式)}  设$\Omega$为平面有界区域,其边界由有限条分段$C^1$的简单闭曲线组成,边界的定向为诱导定向.如果$P,Q$在包含$\overline\Omega$的开集上为$C^1$函数,则
$$\iint_\Omega\left(\frac{\partial Q}{\partial x}-\frac{\partial P}{\partial y}\right)\mathrm{d}x\mathrm{d}y=\oint_{\partial\Omega}P\,\mathrm{d}x+Q\,\mathrm{d}y$$.
\end{ti}


\begin{ti}
\textbf{定理6.23($Gauss$公式)}  设$\Omega$为$\mathbb{R}^3$中的有界区域,其边界由有限个分段$C^1$的闭曲面组成,曲面的定向为诱导定向(即外侧方向).如果$P,Q,R$在包含$\overline\Omega$的开集上为$C^1$函数,则
$$\iiint_\Omega\left(\frac{\partial P}{\partial x}+\frac{\partial Q}{\partial y}+\frac{\partial R}{\partial z}\right)\mathrm{d}x\mathrm{d}y\mathrm{d}z=\iint_{\partial\Omega}P\,\mathrm{d}y\wedge\mathrm{d}z+Q\,\mathrm{d}z\wedge\mathrm{d}x+R\,\mathrm{d}x\wedge\mathrm{d}y$$.
\end{ti}

\begin{ti}
\textbf{定理6.24($Stokes$公式)}  设$\Sigma$为定向曲面,$\Omega$为曲面上的有界区域,其边界赋以诱导定向.如果$P,Q,R$为$\Omega$附近的连续可微函数,则
$$\iint_\Omega\left[\left(\frac{\partial R}{\partial y}-\frac{\partial Q}{\partial z}\right)\mathrm{d}y\wedge\mathrm{d}z+\left(\frac{\partial P}{\partial z}-\frac{\partial R}{\partial x}\right)\mathrm{d}z\wedge\mathrm{d}x+\left(\frac{\partial Q}{\partial x}-\frac{\partial P}{\partial y}\right)\mathrm{d}x\wedge\mathrm{d}y\right]$$
$$=\oint_{\partial\Omega}P\,\mathrm{d}x+Q\,\mathrm{d}y+R\,\mathrm{d}z$$.
\end{ti}


\BookEnd
